\subsection{REAL FORMS}

If \(\mathfrak{a}\) is a complex Lie algebra, we denote by \(\mathfrak{a}_{[\mathbf{R} ]}\) (or sometimes by \(\mathfrak{a}\) ) the real Lie algebra obtained by restriction of scalars. If \(\mathfrak{g}\) is a real Lie algebra, we denote by \(\mathfrak{g}_{(\mathbf{C})}\) (or sometimes by \(\mathfrak{g}_{\mathbf{C}}\) ) the complex Lie algebra \(\mathbf{C}\otimes_{\mathbf{R}}\mathfrak{g}\) obtained by extension of scalars. The homomorphisms of real Lie algebras \(\mathfrak{g}\to \mathfrak{a}_{[\mathbf{R}]}\) correspond bijectively to the homomorphisms of complex Lie algebras \(\mathfrak{g}_{(\mathbf{C})}\to \mathfrak{a}\) : if \(f:\mathfrak{g}\to \mathfrak{a}_{[\mathbf{R}]}\) and \(g:\mathfrak{g}_{(\mathbf{C})}\to \mathfrak{a}\) correspond, we have \(f(x) = g(1\otimes x)\) and \(g(\lambda \otimes x) = \lambda f(x)\) for \(x\in \mathfrak{g},\lambda \in \mathbf{C}\) .

DEFINITION 1. Let \(\mathfrak{a}\) be a complex Lie algebra. A real form of \(\mathfrak{a}\) is a real subalgebra \(\mathfrak{g}\) of \(\mathfrak{a}\) that is an \(\mathbf{R}\) -structure on the \(\mathbf{C}\) -vector space \(\mathfrak{a}\) (Algebra, Chap. II, §8, no. 1, Def. 1).

This means that the homomorphism of complex Lie algebras \(\mathfrak{g}_{(\mathbf{C})}\to\mathfrak{a}\) associated to the canonical injection \(g\to a_{[R]}\) is bijective. Thus, a real subalgebra g of a is a real form of a if and only if the subspaces g and ig of the real vector space a are complementary. The conjugation of a relative to the real form g is the map \(\sigma:a\to a\) such that

\[
\sigma (x + i y) = x - i y, \quad x, y \in \mathfrak {g}.\tag{1}
\]

PROPOSITION 1. a) Let \(\mathfrak{g}\) be a real form of \(\mathfrak{a}\) and \(\sigma\) the conjugation of \(\mathfrak{a}\) relative to \(\mathfrak{g}\) . Then:

\[
\sigma^ {2} = \mathrm{Id} _ {\mathfrak {a}}, \quad \sigma (\lambda x + \mu y) = \bar {\lambda} \sigma (x) + \bar {\mu} \sigma (y), \quad [ \sigma (x), \sigma (y) ] = \sigma [ x, y ]\tag{2}
\]

for \(\lambda, \mu \in \mathbf{C}\) , \(x, y \in \mathfrak{a}\) . An element \(x\) of \(\mathfrak{a}\) belongs to \(\mathfrak{g}\) if and only if \(\sigma(x) = x\) . b) Let \(\sigma: \mathfrak{a} \to \mathfrak{a}\) be a map satisfying (2). Then the set \(\mathfrak{g}\) of fixed points of \(\sigma\) is a real form of \(\mathfrak{a}\) , and \(\sigma\) is the conjugation of \(\mathfrak{a}\) relative to \(\mathfrak{g}\) .

The proof is immediate.

Note that if B denotes the Killing form of a, and if g is a real form of a, the restriction of B to g is the Killing form of g; in particular, B is real-valued on \(g \times g\) . Assume that a is reductive; then the real Lie algebra g is compact if and only if the restriction of B to g is negative (§1, no. 3). In that case we say that g is a compact real form of a.

